\documentclass{article}
\usepackage{amsmath,amssymb}
\usepackage[export]{qam}
\usepackage{qam-physics}
\begin{document}
\qamframe{Lab}
\qamprofile{hc}{Lab}{L_{\rm EM}}{conditional}
  {Real electromagnetic doublet in a fixed laboratory profile with active duality}
\begin{qamcontext}{hc}
\section*{Two-port readout}
\[
\qamformula{ports}{conditional}
  {A=E+B,\quad D=E-B,\quad
   (N_+,N_-,N_E)=(|A|^2,|D|^2,|A+D|^2)}
\]
\[
\qamformula[ports]{field-invariant}{conditional}
  {s=\tfrac12(N_E-N_+-N_-)+\tfrac i2(N_+-N_-)}
\]
\[
\qamformula[field-invariant]{duality-pair}{conditional}
  {(E,B)=(2e_x,e_y),\quad s=3,\quad s_\pm=e^{\pm i\pi/2}s=\pm3i}
\]
\[
\qamformula[duality-pair]{distinct-means}{conditional}
  {\bar s=\tfrac12(s_++s_-)=0,\quad
   M_2=\tfrac12(|s_+|^2+|s_-|^2)=9}
\]
\section*{Weighted zero checks}
Take \(p_j\ge0\) with \(\sum_j p_j=1\).
\[
\qamformula{moments}{conditional}
  {\bar s=\sum_jp_js_j,\quad M_2=\sum_jp_j|s_j|^2,\quad
   M_2-|\bar s|^2=\sum_jp_j|s_j-\bar s|^2}
\]
\[
\qamformula[moments]{branch-zero}{conditional}
  {M_2=0\iff s_j=0\ \text{for every }j\text{ with }p_j>0}
\]
\[
\qamformula[moments]{branch-bound}{conditional}
  {M_2\le\varepsilon^2,\ p_j>0
   \ \Longrightarrow\ |s_j|\le\varepsilon/\sqrt{p_j}}
\]
\[
\qamformula[branch-zero]{nonzero-null}{conditional}
  {\begin{gathered}
     n_j=|E_j|^2+|B_j|^2,\qquad \bar n=\sum_jp_j n_j,\\
     M_2=0,\ \bar n>0
     \ \Longrightarrow\ \exists j:\ p_j>0,\ s_j=0,\ n_j>0 .
   \end{gathered}}
\]
\end{qamcontext}
The final implication certifies existence. For a statement on every
positive-weight branch, require \(n_j>0\) on every such branch. For
example, invertible field transports preserve nonzero input fields.
The profile describes the local field route and its declared readouts.

\begin{qamcontext}{hc}
\section*{Nullity of the coherently reconstructed input}
The two port norms fix the imaginary part of the input invariant.
Cauchy--Schwarz bounds its real part, and every intermediate value is
realized by varying the angle between the port vectors.
\[
\qamformula[field-invariant]{input-range}{conditional}
  {|\operatorname{Re}s|\le\sqrt{N_+N_-},\quad
   \operatorname{Im}s=\tfrac12(N_+-N_-)}
\]
\[
\qamformula[field-invariant]{input-null}{conditional}
  {s=0\iff N_+=N_-\ \text{and}\ N_E=2N_+}
\]
Unequal norms already exclude nullity. If both norms are zero, the
input field vanishes. With equal positive norms, measure the squared
norm of the vector sum or equivalently \(4|E|^2\).
\[
\qamformula[input-null]{input-collision}{conditional}
  {\begin{gathered}
    (N_+,N_-,N_E)=(5,5,16)\ \Longrightarrow\ s=3,\\
    (N_+,N_-,N_E)=(5,5,10)\ \Longrightarrow\ s=0,\quad |E|^2+|B|^2=5 .
   \end{gathered}}
\]
Both inputs have the same projected invariants \(\pm5i/2\) and
the same equally weighted projected second moment \(25/4\).

\section*{Input nullity with bounded measurement errors}
Let \(I_+,I_-,I_E\) be the allowed intervals for the three squared
norms. Assume these intervals give all the error constraints.
\[
\qamformula[input-null]{interval-compatibility}{conditional}
  {\begin{gathered}
    J=I_+\cap I_-\cap(I_E/2)\cap[0,\infty),\\
    \text{A nonzero null input is compatible}
    \iff J\cap(0,\infty)\ne\varnothing .
   \end{gathered}}
\]
Every positive \(q\in J\) is realized by orthogonal electric and
magnetic vectors of length \(\sqrt{q/2}\). Zero field is compatible
exactly when \(0\in J\).
\[
\qamformula[interval-compatibility,input-collision]{interval-examples}{conditional}
  {\begin{gathered}
    (5,5,16)\ \Longrightarrow\ J=\varnothing,\qquad
    (5,5,10)\ \Longrightarrow\ J=[4.95,5.05]
   \end{gathered}}
\]
Here each measured squared norm has absolute error at most \(0.1\).
For the second triple, the invariant is an affine function of the
three errors. Convexity bounds its modulus by the largest value at
the eight corners of the error box.
\[
\qamformula[field-invariant,interval-examples]{sharp-input-bound}{conditional}
  {|s|\le0.15,\qquad
   (N_+,N_-,N_E)=(5.1,5.1,9.9)\ \Longrightarrow\ s=-0.15}
\]
The corner is physically realizable by real port vectors. Thus the
bound is sharp even with the Gram constraints. Compatibility with
exact nullity and a guaranteed upper bound are separate readouts.

\section*{Four norms with shared calibration}
\[
\qamformula[ports,field-invariant]{four-norm-readout}{conditional}
  {\begin{gathered}
    N_B=|A-D|^2,\qquad N_E+N_B=2(N_++N_-),\\
    s=(N_E-N_B)/4+i(N_+-N_-)/2 .
   \end{gathered}}
\]
Take \(\widehat N_j=gN_j+b\) with one common gain \(g>0\)
and one common offset \(b\).
\[
\qamformula[four-norm-readout]{shared-calibration}{conditional}
  {\begin{gathered}
    S_\Delta=(\widehat N_E-\widehat N_B)/4+
      i(\widehat N_+-\widehat N_-)/2=gs,\\
    b=\bigl[2(\widehat N_++\widehat N_-)
           -(\widehat N_E+\widehat N_B)\bigr]/2 .
   \end{gathered}}
\]
The measured tuple \((5.1,5.1,10.1,10.1)\) gives \(S_\Delta=0\)
and \(b=0.1\). Its corrected port norms are positive, certifying
a nonzero null input in this calibration model.

If \(\widehat N_j=gN_j+b+e_j\) with \(|e_j|\le\delta_j\),
the offset still cancels. The remaining error obeys
\[
\qamformula[shared-calibration]{calibration-residual}{conditional}
  {|S_\Delta-gs|\le
  \sqrt{\left(\frac{\delta_E+\delta_B}{4}\right)^2+
        \left(\frac{\delta_++\delta_-}{2}\right)^2}.}
\]
An upper bound on \(|s|\) from these data uses a positive lower
bound on \(g\). The normalized invariant can instead use the field
scale in the same observer profile.
\[
\qamformula[shared-calibration]{calibrated-scale}{conditional}
  {T=\frac{\widehat N_E+\widehat N_B-\widehat N_+-\widehat N_-}{2}
   =gn,\qquad n=|E|^2+|B|^2}
\]
\[
\qamformula[calibrated-scale,shared-calibration]{relative-invariant}{conditional}
  {T>0\ \Longrightarrow\ |S_\Delta|/T=|s|/n}
\]
With residual errors, let \(D_s\) be the bound on \(|S_\Delta-gs|\)
above and \(D_n=(\delta_E+\delta_B+\delta_++\delta_-)/2\).
The same scale readout satisfies \(|T-gn|\le D_n\).
\[
\qamformula[calibration-residual,calibrated-scale]{relative-error-bound}{conditional}
  {T>D_n\ \Longrightarrow\
   \frac{\max(0,|S_\Delta|-D_s)}{T+D_n}
   \le \frac{|s|}{n}
   \le \frac{|S_\Delta|+D_s}{T-D_n}.}
\]
For \((5.1,5.1,10.1,10.1)\) with all \(\delta_j=0.1\), one has
\(T=5\), \(D_n=0.2\), and \(D_s=\sqrt5/20\).
The normalized modulus is therefore at most \(\sqrt5/96\).

\section*{Reference boost from scale and directed flux}
Use \(c=1\). In addition to the norm data, retain the directed
flux \(F_g=2g(E\times B)\) in the same laboratory axes and with the
same gain as the norm readout.
\[
\qamformula[calibrated-scale,shared-calibration]{reference-boost}{conditional}
  {\begin{gathered}
    n^2-|s|^2=4|E\times B|^2,\qquad
    v_{\rm ref}=\frac{F_g}{T+|S_\Delta|}\quad (|S_\Delta|>0),\\
    |F_g|^2=T^2-|S_\Delta|^2 .
   \end{gathered}}
\]
The last equality is a consistency condition on exact flux and
norm data. The positive gain cancels from the boost vector.
\[
\qamformula[reference-boost]{opposite-reference-boosts}{conditional}
  {E=2e_x,\quad B=\pm e_y:\quad
   T=5g,\quad S_\Delta=3g,\quad F_g=\pm4g e_z,\quad
   v_{\rm ref}=\pm e_z/2}
\]
The four norm values agree for these inputs. The directed flux
selects the boost. A nonzero null input instead gives a formal speed
of one and has no finite reference boost.

\section*{Reference boost with a separately calibrated flux}
The flux magnitude is optional for reconstruction. Let \(F_{\rm dir}\)
be any positive multiple of \(E\times B\) in the laboratory axes.
Its gain may differ from that of the norm readout.
\[
\qamformula[relative-invariant,reference-boost]{direction-only-boost}{conditional}
  {\begin{gathered}
     r=|S_\Delta|/T,\qquad
     v_{\rm ref}=\sqrt{\frac{1-r}{1+r}}\,
       \frac{F_{\rm dir}}{|F_{\rm dir}|}\quad (0<r<1),\\
     r=1\quad\Longrightarrow\quad v_{\rm ref}=0 .
   \end{gathered}}
\]
For \(0<r<1\), the flux is nonzero by the field identity above.
The direction must retain its sign. At \(r=1\) no direction
measurement is needed. The magnitude consistency check above still
requires matched gains.
\end{qamcontext}

\section*{Directed flux from opposite observer boosts}
\qamframe{PlusZ}\qamframe{MinusZ}
\qamprofile{plus-z}{PlusZ}{L_{\rm EM}}{conditional}
  {Observer boosted along positive z at speed one quarter}
\qamprofile{minus-z}{MinusZ}{L_{\rm EM}}{conditional}
  {Observer boosted along negative z at speed one quarter}
\begin{qamcontext}{hc}
\[
\qamformula{boost-probe-input}{conditional}
  {E=2e_x,\quad B=b e_y,\quad b\in\{-1,1\},\quad v=1/4}
\]
\end{qamcontext}
\qammap{probe-plus}{Lab}{PlusZ}
\qammap{probe-minus}{Lab}{MinusZ}
\[
\qamphysconvert{plus-scale}{conditional}{probe-plus}{boost-probe-input}
  {plus-z}{n_+=(85-32b)/15}
\]
\[
\qamphysconvert{minus-scale}{conditional}{probe-minus}{boost-probe-input}
  {minus-z}{n_-=(85+32b)/15}
\]
For comparison, record each observer's reported value in the
laboratory record with its observer label. These maps copy the reported
values, preserving their meaning as boosted field scales.
\qammap{report-plus}{PlusZ}{Lab}
\qammap{report-minus}{MinusZ}{Lab}
\[
\qamphysconvert{plus-report}{conditional}{report-plus}{plus-scale}
  {hc}{N_{\rm report,+}=n_+}
\]
\[
\qamphysconvert{minus-report}{conditional}{report-minus}{minus-scale}
  {hc}{N_{\rm report,-}=n_-}
\]
\begin{qamcontext}{hc}
\[
\qamformula[plus-report,minus-report]{probe-flux}{conditional}
  {N_{\rm report,-}-N_{\rm report,+}=\tfrac{64b}{15}
   =8\gamma^2v\, e_z\cdot(E\times B)}
\]
\end{qamcontext}
The x and y probe differences vanish for this pair. The z difference
selects the signed flux direction, giving \(v_{\rm ref}=b e_z/2\).

\section*{Residual flux and a corrective boost}
\qamframe{WrongZ}\qamframe{Rest}
\qamprofile{wrong-z}{WrongZ}{L_{\rm EM}}{conditional}
  {Observer moving along negative z at speed one half relative to the laboratory}
\qamprofile{rest}{Rest}{L_{\rm EM}}{conditional}
  {Flux-free observer reached by the declared corrective boost}
\qammap{wrong-boost}{Lab}{WrongZ}
\[
\qamphysconvert{residual-flux}{conditional}{wrong-boost}{boost-probe-input}
  {wrong-z}{b=1,\quad s=3,\quad E'\times B'=(20/3)e_z}
\]
\begin{qamcontext}{wrong-z}
\[
\qamformula[residual-flux]{corrective-rapidity}{conditional}
  {\eta_{\rm corr}=\tfrac12\operatorname{arsinh}
   \left(2|E'\times B'|/|s|\right),\qquad
   v_{\rm corr}=\tfrac45 e_z}
\]
\end{qamcontext}
\qammap{correct-boost}{WrongZ}{Rest}
\[
\qamphysconvert{rest-flux-check}{conditional}{correct-boost}
  {corrective-rapidity}{rest}{B''=0,\quad |E''|^2=3,\quad E''\times B''=0}
\]
The correction is relative to the intermediate observer. Collinear
velocity addition yields the net laboratory velocity
\[
\frac{-1/2+4/5}{1-(1/2)(4/5)}=\frac12 .
\]
The target flux check here is a calculation. An independent probe
in the target profile would check the realized transformation.

\section*{Identifying the minimum-speed flux-free observer}
In the computed rest profile the electric field points along \(e_x\).
A further boost along this axis preserves the field and its zero
flux, while changing the observer.
\begin{qamcontext}{rest}
\[
\qamformula[rest-flux-check]{flux-free-family}{conditional}
  {\begin{gathered}
     E''=\sqrt3 e_x,\quad B''=0,\quad
     \gamma_{\rm ref}=2/\sqrt3,\\
     \gamma_{\rm target}=\gamma_{\rm ref}\cosh\xi
     \quad\text{for a further field-axis boost of rapidity }\xi .
   \end{gathered}}
\]
\[
\qamformula[flux-free-family]{clock-selection}{conditional}
  {\gamma_{\rm target}=\gamma_{\rm ref}\iff \xi=0}
\]
\[
\qamformula[flux-free-family]{clock-example}{conditional}
  {\tanh\xi=\pm1/4\quad\Longrightarrow\quad
   \gamma_{\rm target}=8/\sqrt{45}>\gamma_{\rm ref}}
\]
\end{qamcontext}
The Lorentz factors compare observer velocities to the laboratory.
Zero field flux tests the flux-free family. The additional velocity
or inertial clock comparison selects its minimum-speed member.
\end{document}
